\section{METHODS}
\label{sec:methods}

\subsection{Design Philosophy}

This research employs Research through Design (RtD) as primary methodological framework. RtD positions the design process itself as inquiry, where knowledge emerges through making and iterating upon designed artifacts \parencite{coulton2017games,zimmerman2007research}. The design process follows \textcite{zimmerman2003play} iterative methodology, treating play as research. Each playtest generates data informing subsequent decisions, creating a spiral of refined understanding.

\subsection{Development Process}

Development unfolded through several phases. In the generation phase, I explore potential mechanics addressing theoretically-identified problems through brainstorming and design conversations with debate and game-design peers. In the critical scrutiny phase, I evaluate whether candidate designs congruently address problems without introducing new difficulties (e.g., judge cognitive load, time limitations), using my own game design praxis awareness and the theoretical framework as evaluation criteria. In the validation phase, surviving elements underwent testing within observed small pilot playtests followed each by a post-session debrief. In the critical analysis phase I synthesize relevant notes on the design based on my observations and opinions expressed during the debrief.

These phases spiraled 3 times (each critical analysis phase sourcing a new generation phase). The first playtest (one debate) was held in March 2026. A second design iteration was drafted based on its findings. The second playtest was conducted in April 2026 with seven veteran debaters in a familiar informal setting, yielding a third design iteration. Both sessions are documented in Annexes A--G.

\subsection{Research Methods}

For analyzing the quality of the design, two playtests were organized. Each took about 4 hours and included a set of workshops preparing players for the debate, the debate itself and the group debrief. During these sessions the research methods used were participant observation and semi-structured post-session group debriefs, followed by qualitative data analysis.

\subsubsection{Participant Observation}

Following \textcite{baker2006observation}, I adopted a participant-as-observer role during both playtest sessions. My researcher status was known to participants, and I participated actively in both sessions to obtain firsthand understanding.

Observation focused on overall debate flow and more granularly on behaviors relevant to entertainment value, rule legibility and adherence ease, and elements distinct from standard formats, judging whether these indicate positive or negative tendencies vis-\`{a}-vis identified problems and the theoretical framework's predictions.

Informed by the advice of \textcite{spradley1980participant}, I recorded field notes in two phases: condensed and expanded. Condensed notes, taken during debates, captured key moments, quotations, and behavioral observations using structured protocols. Expanded notes, written within 24 hours of the sessions, elaborated, added contextual observations, and recorded analytical reflections. Both sessions were audio-recorded with participant consent to enable field note verification and capture verbal content not fully documented in real-time.

\subsubsection{Participants}

Debaters were recruited for the first session with the aid of ARDOR Muntenia, the NGO managing the debate community in Bucharest, and was held as part of a weekend meeting of the Jean Monet highschool debate club. It was held with 8 senior debaters aged 18 to 20. Roles within the debate were assigned randomly. The second session was organized through personal contacts: participants were veteran debaters known to the researcher, recruited informally. The debate was held with 7 senior debaters aged 28 to 38, and roles within the debate were assigned randomly. 

The majority of Romanian debaters are minors. Ideally, the study would involve minor participants, but ethical guidelines pushed toward this strategy as best alternative. In Romania, the young adults debate community is largely homogenous in experience; the sample is thus representative of the format's most likely target audience: experienced, not novice, debaters.

\subsubsection{Post-Session Debriefs}

Semi-structured group debriefs complemented observation. After each play session an about 1 hour long conversation was held with each playtests' participants together: the researcher introduced a set of pre-planned open questions derived from the theoretical predictions (Annex C), but allowed the conversation to develop organically, following participant-introduced threads and encouraging elaboration on emergent themes. Both sessions were audio-recorded with participants' consent.

\subsection{Data Analysis}

Analysis follows a design-oriented comparative approach. The primary purpose is not generating generalizable theory but informing ongoing design decisions and articulating how the designed format produces different experiences than existing formats. This aligns with experimental game design methodology \parencite{waern2015experimental}, where design experimentation becomes research when ``the aim is to understand something generic, some more fundamental aspect of game design.''

Before data collection, the five evaluation fields derived from the research question and theoretical framework---each corresponding to a specific effect the designed format is predicted to produce---were operationalized into a taxonomy of \textit{potential clues}: specific observable behaviors or reported experiences that would count as evidence for or against each field's prediction. Additional clue categories were created after data collection where I observed the taxonomy to not be comprehensive enough. Both positive clues (behaviors consistent with the format achieving its intended effect) and negative clues (behaviors indicating a design gap or failure to achieve it) were specified at this stage. This pre-specification ensures the analysis remains responsive to disconfirming evidence rather than converging only on confirming observations. The complete taxonomy, reproduced in Annex~I, thus defines in advance what would need to be observed to consider the format effective or ineffective in each field. The research subquestion of whether the format achieves its intended effects is answered field by field: the balance of clues found within a field constitutes evidence for or against its prediction, and the aggregate across all five fields answers the whole.

Between playtest sessions, rapid analysis of observation and debrief data focused on actionable insights: what mechanics worked and should be preserved, what did not work and caused confusion or disengagement, and what surprised me and suggested new opportunities. This treats playtesting as ``play as research'' \parencite{zimmerman2003play}, because play emerges from rules as they are inhabited by participants in ways that cannot be fully predicted in advance, each session surfaces questions that were not part of the initial investigation and could not have been anticipated without actually playing. The design cycle is driven not by confirming prior expectations but by discovering what the design still needs to become.

After all playtest sessions were concluded, systematic analysis of field notes and session transcripts proceeded in four steps. First, both sources were read in full and individual clues were extracted: each specific behavior observed or participant statement made that could constitute evidence---positive or negative---for any field's prediction. Second, each clue was assigned to one of the pre-specified taxonomy categories in Annex~I; where the taxonomy proved insufficient, new categories were created. Third, each clue was assessed for evidential strength using explicit criteria: behavioral evidence was weighted above self-report; unsolicited observations above responses to direct prompts; evidence consistent across both playtest sessions above single-session occurrences; and novel or different behaviors above continuities, which may reflect participants' existing training rather than effects of the format \parencite{lankoski2015formal}. Fourth, the clues within each field were synthesized into an argument about whether the field's prediction held, weighting stronger clues more heavily and treating the meaningful absence of predicted negative clues as positive evidence where applicable.

Together, these steps can demonstrate: (a) specific predicted effects of the format did or did not materialize at the behavioral or participant-report level; (b) specific mechanics appear connected to specific observed effects; and (c) observed effects align with the format's intended design goals. They cannot demonstrate persistence, superiority, or generalizability, appropriate limitations for proof-of-concept Research through Design where contribution lies in demonstrating that systematic design can address identified problems and in articulating design principles meriting further investigation or usage.

\subsection{Ethical Considerations}

This research adheres to Uppsala University's ethical guidelines for research involving human participants. All playtest participants were informed of the purpose and context, how data would be gathered, stored, and analyzed, both during and after recruitment. All were able to offer and retract consent, granularly or wholesale, to participation in every process part and to data being collected, stored, or used. Peers whose input was used in developing this format were informed of context and credited. All data is stored securely in encrypted Microsoft OneDrive offered by Uppsala University, separated by content type, and will be deleted once thesis development and presentation is complete. All data has been anonymized and is not publicly available.

To ensure psychological safety, the rigorous ethical standards imposed by the ARDOR Ethical Guide and ARDOR Trainer Safety Guide will be followed. Debate topics will avoid personal trauma topics and debaters will be able to refuse to engage with motions. Debaters will be able to always talk to a third-party equity officer regarding emotional safety issues and report abusive or distressing behavior. During debates, all debaters will be able to exit without explanation.

Beyond these, all participants will have pre-session briefing and post-session debrief providing chances to address potential distressing elements before activity and discuss experience afterwards.

\subsection{Positionality and Reflexivity}

I approach this research as a complete community member \parencite{adler1987membership}. My fifteen years in competitive debate have shaped not only my knowledge but my values, assumptions, and emotional investments. This informs several positions I must acknowledge: insider knowledge enabling nuanced understanding but potentially creating blind spots; reformist motivation emerging from genuine frustration with existing formats, predisposing me toward confirming my critique; a novel disciplinary lens through game design rather than the educational or rhetorical frameworks common in forensics scholarship; and power dynamics created by my status as former Team Romania coach and successful debater, which may influence participant responses. My competitive success within the formats I critique creates additional tension regarding whether my familiarity enables recognition of subtle issues or obscures problems I cannot fully perceive.

To support reflexive awareness, I completed every playtest journal with a reflexive note \parencite{knott2022interviews}, actively seeking data contradicting expectations. The second session introduced a heightened positionality challenge: all participants were close personal acquaintances, which deepened rapport and informal observation access while simultaneously amplifying the risk of social desirability effects and confirmation bias. This proximity is acknowledged as a limitation.

